\section{InferCell: an inference-first design}
\label{sec:design}

\IC{} is written in Julia \citep{bezanson2017julia}. It uses the SciML solvers for
ODEs and jump processes \citep{rackauckas2017differentialequations}. This section
gives the design ideas. The documentation gives the full interface
(\cref{sec:availability}). \Cref{fig:architecture} shows the parts and how they
connect.

\begin{figure*}[t]
\centering
\placeholderfig{0.9\textwidth}{6cm}{Figure 1 (F11): the architecture. Parameters
with provenance feed modules; modules declare formalism and edges; the resolver
checks the composition; the composed model runs as the exact hybrid or as one
mean-field ODE; inference routines attach to each data stream. To be drawn.}
\caption{\textbf{The \IC{} architecture.} \placeholder{Caption to be written with the
figure.}}
\label{fig:architecture}
\end{figure*}

\subsection{Parameters are distributions with a source}

Each parameter in \IC{} is an object that holds a value, a prior distribution, a
role and a record of its source. The record names the file, the table and the
identifier the value came from, and the values that other files give for the same
identifier. Each parameter also has one of four labels for how well its prior is
informed:
\begin{itemize}
  \item \emph{balanced}: the value and width come from a balanced parameter table;
  \item \emph{prior default}: the source gives no width, so the prior is wide;
  \item \emph{asserted}: the source gives no uncertainty, so we set the prior;
  \item \emph{not imported}: the parameter is not in the source.
\end{itemize}
The loader does not choose silently between two files that disagree. It stops and
asks for a decision, and the decision is then part of the record. A user can ask
the composed model which priors no measurement supports. That question has a
direct answer, because the answer is in the model and not in a supplement.

\subsection{Modules couple through typed edges}

A module declares its states, its parameters and its formalism: ODE or jump. It
also declares how it touches the rest of the cell. \IC{} has \pn{fw.edgekinds} kinds
of coupling edge (\cref{tab:edges}). Each kind says what crosses the edge and on
which clock. Mass and currency edges act inside the ODE solve. Catalytic edges
and deferred costs cross every second, at the handshake between the two blocks.
Rate-constant edges cross every minute, when the stochastic block rebuilds its rate
constants from the metabolite pools. The set of kinds is closed. A module cannot
invent a new kind of coupling without a change to the framework.

\begin{table}[t]
\caption{The \pn{fw.edgekinds} edge kinds and their clocks.}
\label{tab:edges}
\centering
\small
\begin{tabular}{@{}p{0.27\columnwidth}p{0.45\columnwidth}p{0.18\columnwidth}@{}}
\toprule
Edge & What crosses & Clock \\
\midrule
Mass & a shared continuous state & ODE solve \\
Currency & flux through a shared pool (ATP, GTP) & ODE solve \\
Catalytic & a protein count into a rate law & 1\,s \\
Deferred cost & jump-side cost, debited one step later & 1\,s \\
Rate constant & pools into stochastic rate constants & 60\,s \\
Volume & membrane proteins into cell geometry, and back & each step \\
Clamped & a dependence replaced by a constant & declared only \\
\bottomrule
\end{tabular}
\end{table}

Before anything runs, a resolver reads all the edges and builds the coupling graph.
It checks the composition for about \pn{fw.resolvererrors} kinds of error. Examples are a state that
two modules integrate, an edge that names a species no module owns, and a moiety
that is consumed but never produced. In complete mode it refuses to build a model
with an unowned state or a stranded moiety. A composition that runs is therefore
a composition that the framework has checked.

\subsection{Every reduction is declared}

Any model of a cell leaves things out. In \IC{} each departure from the published
model is an object, not a comment. A reduction label has a category (one of
\pn{fw.redcats}, such as a clamp, a lumped step, a smoothed counter or an asserted
prior) and a reason. The framework collects labels from three places: the edges,
the parameter records and the modules. One call returns the full list for a
composed model. This list is how we report the cost of each choice in
\cref{sec:results}.

\subsection{One interface, two descriptions}
\label{sec:twodesc}

Gene expression in \IC{} has one interface and two descriptions.

\paragraph{The exact hybrid.} The stochastic modules run as a jump process, and the
metabolic modules run as ODEs. A driver steps the two blocks in turn, as the
published model does. Every second it writes protein counts into the metabolic rate
laws, steps the ODEs, debits the costs of the stochastic events, and steps the jump
process. Every minute it rebuilds the stochastic rate constants from the pools. We
did not wrap this in one solver object on purpose. A single jump problem over an ODE
problem would let the propensities follow the pools continuously. That is a
different model from the published one. The driver can also record a stochastic path
and replay it exactly. A replay matches the original run bit for bit at every
handshake. The exact sampler (\cref{sec:sampler}) depends on this.

\paragraph{The mean-field swap.} The gene-expression network of the model is first
order. So the equations for the means of the transcript and protein counts close
exactly, as linear ODEs. \IC{} derives these equations from the same propensities
and update rules that the jump modules use, so the two descriptions cannot drift
apart. With the swap, the whole cell becomes one ODE problem with callbacks for the
one-minute rebuild. Forward-mode automatic differentiation
\citep{revels2016forwardmode} passes through this problem and its callbacks. We test
the gradients against central differences to a relative error of $\pn{fw.adtol}$.

\paragraph{What is and is not differentiable.} The exact hybrid is not
differentiable, because the published model is not. The published model clips a
drained pool at zero and rounds pools to whole molecules. Both make the model a step
function of its parameters. The exact sampler therefore uses no gradients. The
mean-field description gives gradients for the parts of inference that need them.
The cost of the swap is a declared reduction, and we measure it
(\cref{sec:jensen}).

\subsection{Inference attaches to each data stream}
\label{sec:api}

A data stream declares its observation model. In this paper there are two streams:
exact transcript counts per cell, and bulk metabolite measurements of the mean over
many cells. \IC{} has two inference routes.

\begin{itemize}
  \item \textbf{The exact route.} A three-block Gibbs sampler conditions on the
  stochastic path. It updates the paths by conditional sequential Monte Carlo
  \citep{andrieu2010particle}, the metabolic constants by slice sampling, and the
  gene-expression constants from exact one-dimensional conditionals
  (\cref{sec:sampler}).
  \item \textbf{The approximate route} (in development). The paths are integrated
  out. The transcripts keep an exact likelihood, and the bulk metabolites get a
  Gaussian likelihood centred on the mean-field model. The route is designed to use
  the gradients of the swap with the No-U-Turn Sampler \citep{hoffman2014nouturn}.
  \placeholder{B5: built and tested; give the interface in one sentence.}
\end{itemize}

\placeholder{B6, the dispatcher: a module declares its formalism and a stream
declares its observation model; the dispatcher picks the route. Not built. If it is
not built before submission, reword this section as a design principle and remove
the claim.}
